\section{系统概述}
\label{sec:system_overview}

\subsection{系统目标}
智能城市生成系统的设计目标不是单纯提供一个建模界面，而是构建一个能够支撑“意图输入 - 任务转换 - 三维执行 - 结果反馈 - 版本沉淀”的软件系统。结合课程主题，系统需要同时满足以下目标：
\begin{enumerate}[label=\arabic*.]
    \item 面向建模师提供高效率的基础场景生成与局部修改能力。
    \item 面向分析师提供清晰的结果查看、仿真比较与导出能力。
    \item 面向管理员提供权限、资源、模型服务与日志治理能力。
    \item 通过结构化任务机制连接上层交互和底层 Blender 执行。
    \item 通过 LLM 降低表达门槛，但保留确认、纠错和回退能力。
\end{enumerate}

\subsection{用户视角概述}
系统的三类核心角色分别具有不同的使用重心。场景建模师关注生成效率、参数控制和多轮修改；行业分析师关注展示效果、方案比较和导出结果；系统管理员关注权限隔离、资源状态、服务可用性与审计追踪。因此，系统在设计上不能把三类用户混成统一入口下的同质操作流，而应围绕角色工作台进行组织。

\begin{figure}[t]
\centering
\includegraphics[width=\textwidth]{figures/1.png}
\caption{系统整体业务概览}
\label{fig:system_overview}
\end{figure}

\subsection{核心设计思想}
本系统的软件设计遵循以下几项核心思想。
\begin{enumerate}[label=\arabic*.]
    \item 角色驱动：先按角色划分工作台，再在工作台内部组织操作链路。
    \item 任务驱动：所有核心能力最终都归约为可跟踪的任务，而不是零散按钮操作。
    \item 执行解耦：主系统负责理解和编排，Blender 插件负责几何与场景执行。
    \item 数据留痕：项目、场景、任务、日志和解析结果都需要保留版本和过程信息。
    \item 智能受控：LLM 负责辅助理解，不直接越过确认机制修改场景。
\end{enumerate}

\subsection{运行上下文}
从运行形态上看，智能城市生成系统由五类要素共同组成：用户界面层、任务组织层、三维执行层、数据资源层和智能解析层。用户通过主系统表达设计意图；主系统将输入转换为结构化任务；Blender 插件依据任务完成场景生成与修改；资产库、项目库和日志库负责支撑数据持久化；LLM 服务提供自然语言与多模态输入的解释能力；仿真模块在执行结果基础上补充动态表现。
